\documentclass[onecolumn]{article} %Este comando fabrica por defecto el tipo de documento, en este caso un artículo, en el formato de dos columnas
\usepackage[spanish]{babel}
\usepackage[left=2cm, right=2cm, top=2cm, bottom=2cm]{geometry} %Esto es para ponerle los márgenes, puedes hacerlo más grande o más chico
\usepackage{fancyhdr} %no borres esto, es para estilizar el encabezado y el pie de página
\usepackage{graphicx} %este es para adjuntar imágenes
\usepackage{float}
\usepackage{subfig} 
\usepackage{amssymb}
\usepackage{gensymb}
\usepackage{svg}
\usepackage[colorlinks=true, linkcolor=blue, urlcolor=blue, citecolor=blue]{hyperref}

% Configuración del encabezado solo en la primera página
\fancypagestyle{firstpage}{
	\pagestyle{fancy} 
	\fancyhead{}
	\fancyhead[L]{\textbf{Universidad de Antioquia}} % Fecha en la esquina izquierda
	\fancyhead[R]{\textbf{\today}} % Nombre de la institución en la derecha
	\fancyfoot[C]{\thepage} % Número de página centrado
}

\begin{document}
	
	\title{
		\textbf{Mapa de procesos y estructura organizacional}\\ \textbf{Seguros SURA Colombia (Suramericana S.A.)}
		}
		
	
	\author{
		\centering
		\begin{minipage}{0.45\textwidth}
			\centering
			\textbf{David Alejandro Morón Acacio}\\
			Facultad de ingeniería\\
			\texttt{david.moron@udea.edu.co}\\
			\vspace{10pt}  
			\textbf{Luisa María Bohórquez Ardila}\\
			Facultad de ingeniería\\
			\texttt{luisam.bohorqueza@udea.edu.co}\\
		\end{minipage}%
		\hfill
		\begin{minipage}{0.45\textwidth}
			\centering
			\textbf{Iván René Acosta Lallemand}\\
			Facultad de ingeniería\\
			\texttt{ivan.acosta@udea.edu.co}\\
			\vspace{10pt}
			\textbf{Juan David Moreno Díaz}\\
			Facultad de ingeniería\\
			\texttt{juan.morenod@udea.edu.co}
		\end{minipage}
	}
	
	\date{} %este comando es necesario para borrar la fecha que maketitle pone por defecto
	
	\maketitle
	
	\thispagestyle{firstpage}
	
	\section{Desarrollo}
	
	Para este trabajo elegimos a la empresa Seguros SURA Colombia (Suramericana S.A.) la cual se dedica a la prestación de servicios de seguros, gestión de tendencias y riesgos, seguridad social y salud. Esta empresa fue escogida debido a que es una organización líder en Colombia con un nivel de madurez corporativa muy alta, lo cual permite un buen análisis administrativo que exige el curso de habilidades gerenciales. Su tamaño y complejidad permiten observar claramente cómo se dividen las cargas de trabajo y cómo operan los procesos estratégicos, misionales y de apoyo en un entorno real.\\
	
	\subsection{Organigrama}
	
	\begin{figure}[H]
		\centering
		\includesvg[width=1\columnwidth]{Organigrama}
		\caption{Organigrama - Construcción propia con Lucid}
		\label{fig:Og}
	\end{figure}
	
	
	
	Por motivos prácticos se decidió hacer un análisis sólo de la alta gerencia de SURA y de los principales encargados del sector financiero. En la imagen \ref{fig:Og} se muestra el diagrama organizativo que se analiza en este informe. Información extraída de \cite{mp5}, \cite{mp6}. 
	
	\subsection{Tipo de estructura organizacional}
	
	Consideramos que SURA tiene una estructura formal, está compuesta por las personas que conforman la empresa y su relación entre sí. En SURA existen los principios organizacionales de actividades, divisón de trabajo, autoridad, responsabilidad, delegación, unidad de mando y jerarquía, algunas de estas cualidades son observables en el organigrama anteriormente expuesto, por ejemplo, es notoria la jerarquía porque hay una persona que es el principal resposable y a su vez las personas a su cargo que también son responsables de otras personas, en ese orden de ideas, notamos que existe la unidad de mando y la división de trabajo porque según los cargos, se realizan tareas específicas y siempre hay un encargado de supervisar las tareas de cada grupo. \\
	
	Una vez visto el organigrama que se analiza en este informe, es evidente que la estructura organizacional observable es funcional, debido a que sus labores están divididas en unidades muy especializadas, cada una al mando de un jefe independiente que coordina a su equipo de trabajo. Esto es beneficioso para la empresa debido a que cada grupo de trabajo se concentra exclusivamente en sus tareas, aumentando la eficiencia y productividad de los trabajadores, así como una alta especialización en cada producto (salud vs autos). Otro beneficio es que permite una mejor supervisión técnica y favorece la comunicación directa de los empleados con sus respectivos jefes, evitando situaciones incómodas como enterarse de las decisiones de la empresa por medio de ``chismes de pasillo'', también disminuye la presión sobre el jefe al compartir las responsabilidades. \\
	
	A pesar de que no es evidente en las líneas del organigrama, se supone que las unidades deben trabajar en conjunto, seguramente el gerente de planeación financiera trabajará en conjunto algunas veces con la gerente de inversiones y tesorería, así como la vicepresidenta de riesgos estará relacionada con la gerente de asuntos legales. Teniendo eso en cuenta podrían relacionarse en el orgnaigrama de la empresa a los jefes de trabajo afines para que se note el tipo de estructura matricial porque en la práctica algunas áreas son transversales, como las de talento humano, tecnología y finanzas, por lo que no podría ser considerada una estructura funcional como lo indica el organigrama de la empresa. \\
	
	Debido al tamaño de SURA es coherente que elijan un tipo de estructura ``funcional'' gracias al gran tamaño de la organización y su gran número de trabajadores, sin embargo este modelo tiene varios aspectos a mejorar. Uno sería modificar la comunicación entre divisiones debido a que en esta estructura pueden ocurrir silos de información y la burocracia puede retrasar la toma de decisiones ágiles. También pueden ocurrir conflictos en las funciones generales de la empresa, o se podría potenciar la competencia y la rivalidad, provocando así tensiones en el ambiente laboral.
	
	\subsection{Mapa de procesos}
	
	\begin{figure}[H]
		\centering
		\includesvg[width=1\columnwidth]{MapaDeProcesos}
		\caption{Mapa de procesos - construcción propia}
	\end{figure}
	
	Este mapa de procesos fue realizado con información extraída de fuentes directas de la empresa que hablan de los procesos de sostenibilidad\cite{mp1}, que muestran el desempeño en infraestructura y control de la información digital\cite{mp2}, los cuales dan cuenta de sus principios y planeación\cite{mp3}, y de los informes que justifican y describen los procesos que se llevan a cabo en pro de lograr los objetivos empresariales\cite{mp4}.\\
	
	Se identifica que en esta empresa las entradas eran las tendencias y condiciones del entorno, el marco regulatorio, el capital de los usuarios, los recursos tecnológicos y los insumos médicos debido a que estos factores no dependen de la empresa pero sí condicionan el normal funcionamiento de la misma, es fundamental la necesidad del servicio de salud de parte de los clientes, que estos tengan el recurso para acceder a los servicios, que existan insumos médicos con qué atender a esos pacientes y tener en cuenta para cualquier decisión el marco que regule a la institución.\\
	
	Para el caso de SURA la satisfacción del cliente dependía completamente del bienestar y la seguridad para el usuario, de la atención efectiva, de la sostenibilidad y el impacto social postivo, mientras que para los accionistas se debía cumplir la expectativa de mantener la sostenibilidad económica, por lo que identificamos estas como nuestras salidas.\\
	
	Las actividades de largo alcance que se identificaron y que tienen una relación directa con la razón de ser de la empresa y los aspectos que la diferencian de la competencia son: el direccionamiento estratégico, la gestión de tendencias y riegos, la innovación corporativa, la gestión de sostenibilidad ambiental y la gestión de marca y ética. Nótese que todos estos son procesos estratégicos debido a que se planifican al mediano y largo plazo y están íntimamente relacionados con la estructura organizacional, las formas de financiación, las alianzas estratégicas y las políticas organizacionales.\\
	
	Por otro lado, identificamos que el diseño de soluciones y productos, la suscripción y afiliación, la gestión de siniestros y reclamos, la gestión y prestación de salud, y la afiliación, fidelización y comercialización son procesos misionales porque están centrados en generar valor agregado a los clientes, para los usuarios es claro que SURA, entre mcuhas otras cosas, es una aseguradora que se encarga de la prestación de servicios de salud, por lo que las actividades antes mencionadas están relacionadas con el funcionamiento normal de la empresa acorde a la prestación de servicios.\\
	
	Por último, identificamos como procesos de apoyo: la gestión de talento humano, la gestión de tecnología, la transformación digital y ciberseguridad, la gestión de abastecimiento, la gestión jurídica y a la gestión financiera. Las anteriores actividades están presentes para que los procesos estratégicos y misionales puedan funcionar de manera adecuada, aquí están presentes esas áreas transversales que se mencionaron en el análisis del organigrama.
	
	
	\subsection{Estilo de administración}
	
	Acorde a lo visto hasta el momento en el curso de habilidades gerenciales, hay dos estilos de administración que son compatibles con SURA: la teoría humanista y la burocrática. Sura trata a sus trabajadores con dignidad humana y los reconoce como hombre y mujeres sociales, este acto es coherente con la teoría humanista que buscaba nuetralizar la deshumanización del trabajo. Si bien se busca la eficiencia, no es necesario explotar a las personas que conforman a la empresa. Asimismo la empresa posee una organización jerárquica donde hay unidades especializadas en sus respectivos campos, lo cual encaja con la teoría burocrática.\\
	
	También es notorio que SURA utiliza una administración participativa y estratégica, se basan en el empoderamiento de los colaboradores, un enfoque humanista (prioriza el talento humano y el bienestar) y el capitalismo consciente. Una gran herramienta es su enfoque en el desarrollo humano porque genera un alto sentido de pertenencia y baja rotación laboral. Sin embargo, al ser participativo y buscar consensos, la implementación de cambios operativos drásticos puede tomar más tiempo del esperado frente a competidores más agresivos, lo cuál no es algo beneficioso para la empresa.
	
	\subsection{Gerencia estratégica}
	
	\subsubsection{Misión}
	Mejoramos la calidad de vida, a través de la conservación y el mejoramiento de la salud.
	
	\subsubsection{Visión}
	Somos un grupo financiero enfocado en construir y desarrollar un portafolio balanceado con visión de largo plazo, cuyo eje principal son los servicios financieros. Nuestro propósito es crear bienestar y desarrollo armónico para las personas, las organizaciones y la sociedad.
	
	\subsubsection{Principios corporativos (\cite{mp7}, \cite{mp8})}
	\begin{itemize}
		\item \textbf{Equidad:} Base para establecer relaciones ganar-ganar.
		\item \textbf{Transparencia:} Criterio para manejar la información, tener la disposición a ser examinados y ser claros sobre la forma en que gestionamos los negocios.
		\item \textbf{Respeto:} Define la manera cómo entendemos y valoramos al otro en toda su dignidad.
		\item \textbf{Responsabilidad:} Talento humano con las competencias para desarrollar la estrategia con identidad cultural de Grupo SURA.
	\end{itemize}
	
	\subsubsection{Valores corporativos}
	\begin{itemize}
		\item Respeto por sí mismo y por los grupos de personas con los que interactuamos.
		\item Ética como fundamento de nuestras actividades.
		\item Actitud permanente de servicio.
		\item Trabajo en equipo.
		\item Compromiso hacia el logro de resultados.
		\item Coherencia entre lo que pensamos, decimos y hacemos.
	\end{itemize}
	
	
	\subsection{Planeación estratégica}
	
	Sura no solo vende seguros, su planeación estratégica se centra en la "Gestión de Tendencias y Riesgos". Buscan anticiparse a los problemas de los clientes (salud, movilidad, competitividad) en lugar de solo pagar cuando ocurre el accidente. Su visión a largo plazo está centrada en la sostenibilidad y la transformación digital. Son pioneros en observar el entorno (megatendencias globales). Alinear a miles de empleados a nivel nacional bajo esta misma planeación requiere un esfuerzo titánico de comunicación interna. \\
	
	En este orden de ideas creemos que la planeación estratégica que puede estar utilizando SURA es el balanced scorecard, debido a que SURA se centra en los cuatro aspectos de este tipo de planeación presentada en clase: primero, su foco son los clientes, el cual es un punto de esta planeación; segundo, busca preveer daños, lo cual mejora las finanzas de la empresa, siendo las finanzas el segundo aspecto; tercero, como su visión estratégica busca la continua mejora, SURA se encuentra en constante proceso de transformación (gestión del talento humano, los sistemas de información y el conocimiento) lo cual coincide con el punto del aprendizaje y cremiento de la empresa; el cuarto aspecto es el proceso interno, el cual es notorio en esta empresa porque son pioneros en observar el entorno y a partir de ahí sacar nuevos productos y servicios.\\
	
	Con esta planeación estratégica SURA tiene en claro sus oportunidades porque tiene el foco puesto en el cliente sin descuidar el proceso interno, sin embargo, en estipo de planeación estratégica no se tiene en cuenta la brecha operativa y la brecha estratégica que seguramente exista entre las expectativas de la empresa y el resultado de los procesos llevados a cabo. Podrían adicionar un análisis GAP (es decir, de brechas) para complementar su proceso de planeación.
	
	\subsection{Sugerencias}
	
	A pesar de que ya se han realizado críticas a lo largo del documento se condensan las críticas y se añaden las sugerencias en términos de estructura y estilo de administración  en este apartado: 
	
	\begin{itemize}
		\item Se considera que debe hacer una modificación al organigrama porque coincide con una estructura organizacional funcional pero a nivel prático llega a ser matricial.
		\item Podrían tener una persona encargada o un equipo de trabajo transversal que se encargue únicamente de la comunicaciones entre divisiones para evitar los silos de información.
		\item Podrían contratar a más personas en los procesos administrativos que sean más lentos para que no les perjudique tanto la burocracia.
		\item Mantener un seguimiento de la salud mental de los trabajadores para supervisar y prevenir que no haya un mal ambiente laboral o contratar un psicólogo que pueda mediar los conflictos cuando se presenten.
		\item Implementar ``Células Ágiles'' interdisciplinarias para el lanzamiento de nuevos productos, reduciendo los niveles jerárquicos de aprobación.
		\item Integrar herramientas de Inteligencia Artificial que automaticen tareas operativas rutinarias, permitiendo que el modelo de administración participativa se enfoque 100\% en la creatividad y solución de problemas complejos del cliente.
		\item Complementar la planeación estratégica con un análisis GAP para hacer una comparación clara entre el desempeño real y el planeado por la empresa.
		
	\end{itemize}
	
	
	\subsection{Trabajo de equipo en la actividad}
	
	El anterior líder le preguntó a Iván si podía ser el líder para la actividad y él asumió la responsabilidad sin problema. El líder se encargó de sugerir la empresa que terminó siendo la elegida. Iván ha sido un líder humanista y con iniciativa al momento de empezar la actividad, luego utilizó la estructura funcional, se delegaron tareas y se realizaron dos llamados de control a lo largo del trabajo para verificar cómo iban los avances en la actividad.\\
	
	Se aprovecha este punto para aclarar que \textbf{Los Caribes} sí se identifican con SURA porque su enfoque no es solo transaccional, sino que valora el bienestar de la sociedad y sus empleados, valores que compartimos como futuros profesionales. Asimismo maneja un sistema burocrático que permite la delegación de tareas y especialización de campos, lo cual les parece lo más óptimo para una empresa con muchos trabajadores. A pesar de que algunos procesos se demoren más de lo deseado, al parecer del grupo, es necesario manejar ese control, pero se puede mejorar con mayor personal en las áreas dónde se encuentren los retrasos. 
	
	
	\newpage
	
	\begin{thebibliography}{99}
		
		\bibitem{mp1} https://www.gruposura.com/sostenibilidad/
		
		\bibitem{mp2} https://www.epssura.com/
		
		\bibitem{mp3} https://www.gruposura.com/nuestra-compania/etica-y-gobierno-corporativo/
		
		\bibitem{mp4} https://www.gruposura.com/relacion-con-inversionistas/
		
		\bibitem{mp5} https://www.sura.co/somos-sura/nosotros/gobierno-corporativo
		
		\bibitem{mp6} https://www.sura.co/somos-sura/estructura-organizacional
		
		\bibitem{mp7} https://seguros.sura.cl/compania/principios
		
		\bibitem{mp8} https://www.epssura.com/embarazo-menuparalasmamas-229/preguntas-frecuentes-durante-el-embarazo-menuparalasmamas-233?task=section\&id=29
		
		\bibitem{mp9} https://www.epssura.com/valores-corporativos-menu-boletin2-243
		
	\end{thebibliography}
	
	
\end{document}